\documentclass[tikz,border=6pt]{standalone}
\usepackage{peabra}
\begin{document}
\begin{tikzpicture}[x=1cm,y=1cm,
  k/.style={doboz=32mm, font=\footnotesize\bfseries, minimum height=10mm},
  l/.style={szurkedoboz=32mm, font=\scriptsize, minimum height=14mm, align=center}]
  \node[sotetdoboz=46mm, font=\footnotesize\bfseries] (u) at (0,2.6) {nem biztonságos cselekvések};
  \node[k] (sz) at (-4.2,1.0) {szándékolatlan};
  \node[k] (sd) at (4.2,1.0) {szándékolt};
  \draw[nyilv] (u.south) -- (sz.north); \draw[nyilv] (u.south) -- (sd.north);
  \node[l] (m) at (-6.3,-1.2) {\textbf{megcsúszás}\\figyelmi hiba: jó szándék,\\rossz végrehajtás};
  \node[l] (k1) at (-2.1,-1.2) {\textbf{kihagyás}\\emlékezeti hiba: egy lépés\\kimarad, elfelejtődik};
  \node[l] (t) at (2.1,-1.2) {\textbf{tévedés}\\rossz terv: szabály- vagy\\tudásalapú hiba};
  \node[l] (s) at (6.3,-1.2) {\textbf{szabályszegés}\\szándékos eltérés:\\rutin vagy kivételes};
  \draw[nyilv] (sz.south) -- (m.north); \draw[nyilv] (sz.south) -- (k1.north);
  \draw[nyilv] (sd.south) -- (t.north); \draw[nyilv] (sd.south) -- (s.north);
  \draw[peKiemel, line width=1pt, rounded corners=3pt] (-8.05,-2.15) rectangle (-0.35,-2.75);
  \node[font=\scriptsize, text=peKiemel] at (-4.2,-2.45) {készségalapú szint (Rasmussen)};
  \draw[peKiemel, line width=1pt, rounded corners=3pt] (0.35,-2.15) rectangle (3.85,-2.75);
  \node[font=\scriptsize, text=peKiemel] at (2.1,-2.45) {szabály- és tudásalapú szint};
  \node[felirat, text width=150mm, align=center] at (0,-3.35) {A hibák (megcsúszás, kihagyás, tévedés) és a szabályszegések eltérő okokból fakadnak, ezért eltérő ellenszert igényelnek: tervezést, emlékeztetőt, képzést vagy szervezeti változtatást.};
\end{tikzpicture}
\end{document}
